\chapter{Discussão e Limitações}
\label{chap:discussao}

\orientacao{Capítulo curto e obrigatório. Responde a três perguntas: o que o resultado significa, sob quais condições vale e o que ainda não está demonstrado.

Escreva cada limitação com a causa técnica e o efeito sobre o resultado, em parágrafo próprio com título em negrito. Registre o fato, sem culpado e sem pedido de desculpa: ``o fornecimento do hardware ocorreu em junho, o que concentrou os ensaios em duas semanas'' é registro; ``a demora do fornecedor inviabilizou o trabalho'' atribui culpa.

Limitação declarada com sobriedade aumenta a credibilidade do documento e delimita por escrito o que não foi prometido. Limitação omitida reaparece na reunião de aceite, quando admiti-la custa muito mais.}

\section{Discussão}

\ph{Interpretação integrada dos resultados dos capítulos técnicos, depois de os dados terem sido apresentados.}

\section{Limitações}

\paragraph{\ph{Nome da limitação}.} \ph{Causa técnica da limitação e efeito sobre o resultado.}

\paragraph{\ph{Nome da limitação}.} \ph{Causa técnica da limitação e efeito sobre o resultado.}
